\cvheader{Arsenij Schützer}{Fullstack-Entwickler}{%
  \begin{tabular}{@{}l@{\hspace{0.6em}}l@{}}
    {\color{dim}Adresse:} & Königsberger Str.~4B, 75181 Pforzheim\\
    {\color{dim}Geboren:} & 08.10.1992 in Omsk (Russland)\\
    {\color{dim}Telefon:} & +49 162 629 1630\\
    {\color{dim}E-Mail:}  & \href{mailto:senyaarseniymail@gmail.com}{senyaarseniymail@gmail.com}\\
    {\color{dim}Web:}     & \href{https://senyaak.work}{senyaak.work} \textcolor{dim}{·} \href{https://github.com/senyaak}{github.com/senyaak} \textcolor{dim}{·} \href{https://linkedin.com/in/senyaak}{linkedin.com/in/senyaak}
  \end{tabular}
}

\section*{Kurzprofil}

Fullstack-Entwickler mit 8 Jahren Berufserfahrung. Der Einstieg ins Fach erfolgte über ein eigeninitiativ begonnenes Fernstudium (ils.de, 2014/15) und eine dreijährige Ausbildung zum Fachinformatiker für Anwendungsentwicklung (2015–2018). Schwerpunkte: Angular, Vue, TypeScript, Node.js, Java, PostgreSQL. Erfahrung in der eigenverantwortlichen Umsetzung kompletter Anwendungen (Frontend, Backend, Datenbank und Architektur) sowie in Cross-Functional Teams. Russisch als Muttersprache.

\section*{Technische Kenntnisse}

\cvkvblock{
\cvkv{Hauptstack}{Angular (alle Versionen seit AngularJS), Vue~3, TypeScript/JavaScript, Node.js, Java, PostgreSQL, MongoDB}
\cvkv{Erfahren}{React, NestJS, Python (FastAPI, Kedro), Linux, Git, Docker, WebSocket, REST APIs, JSON Schema}
\cvkv{Weitere Erfahrung}{C, C++, C\#, PHP, Rust, Unity, Jenkins, Kubernetes, Socket.IO, LLM Agents}
}

\section*{Berufserfahrung}

\cvjob{e.telligent GmbH, Regensburg}{09/2023 – heute}{Fullstack-Entwickler}{e.telligent entwickelt Softwarelösungen für den Bereich elektrifizierte Mobilität. Arbeit am Produkt e.CoSys, einer ASPICE-konformen Application-Lifecycle-Management-Plattform für die Automobilentwicklung mit durchgängiger Traceability entlang des V-Modells.}
\begin{itemize}
  \item Fullstack-Entwicklung über den gesamten Technologie-Stack: \textbf{Vue~3 + TypeScript} im Frontend, \textbf{Java} mit \textbf{PostgreSQL} im Backend
  \item Konzeption und Umsetzung kundenspezifischer Features auf einer generischen, lizenzierbaren Plattform-Architektur
  \item Eigenverantwortliche Entwicklung komplexer Frontend-Komponenten
  \item Backend-Entwicklung mit Schwerpunkt auf komplexen Datenstrukturen
  \item Code Reviews und Refinement-Sessions im Team
  \item KI-Service-Wrapper in Python
\end{itemize}

\cvjob{Mittwald CM Service GmbH \& Co.~KG, Espelkamp}{07/2022 – 12/2022}{Fullstack-Entwickler}{Mittwald ist ein etablierter deutscher Webhosting-Anbieter. Mitarbeit am Customer-Facing-Portal als Teil eines größeren Entwicklungsteams.}
\begin{itemize}
  \item Frontend-Entwicklung am Kundenportal mit \textbf{React} und \textbf{TypeScript}
  \item Mitarbeit in einem agilen Team mit etablierten Code-Standards und CI/CD-Prozessen
\end{itemize}

\cvjob{saltation GmbH \& Co.~KG, Bielefeld}{08/2018 – 01/2022}{Fullstack-Entwickler}{saltation entwickelt individuelle Softwarelösungen für verschiedene Branchen (20–40 Mitarbeiter). Eigenverantwortliche Umsetzung mehrerer Projekte vom Frontend bis zur Datenbank, in einigen Fällen Tandemarbeit (Solo-Frontend mit Solo-Backender).}
\begin{itemize}
  \item Content-Management-System für Objektdatenpflege (Häfen, Fahrräder): generische Modelle aus JSON-Schema mit automatischer Formulargenerierung über JSON-Form-Konfiguration. Stack: \textbf{AngularJS, Node.js, PostgreSQL}; vollständige Solo-Verantwortung über mehrere Jahre
  \item Architekturelle Erweiterung des CMS um Internationalisierung (mehrsprachige Inhalte parallel): Anpassungen an Datenbank, Backend und Frontend in koordinierter Umsetzung
  \item System zur Modellierung von Geschäftsprozessen (Schulungen, Weiterbildung): \textbf{AngularJS} mit individueller Engine, mehrere Frontend-Varianten
  \item Administrations-Tool für Geräteverleih: \textbf{Angular} (ab Version~2)
  \item Onboarding-App für Tablets in \textbf{Unity (C\#)}
  \item Gelegentliche Backend-Arbeit in \textbf{Rust} (serverseitige Anpassungen) sowie kleinere Projekte in Java und C\#
\end{itemize}

\cvjob{saltation GmbH \& Co.~KG, Bielefeld}{08/2015 – 07/2018}{Ausbildung: Fachinformatiker für Anwendungsentwicklung (IHK)}{}
\begin{itemize}
  \item Abschlussprojekt: Konzeption und vollständige Umsetzung der Internationalisierungs-Erweiterung für ein bestehendes CMS (Datenbank, Backend, Frontend; Solo-Projekt)
  \item Praxiserfahrung in mehreren Web-Projekten mit AngularJS, Node.js, PostgreSQL
\end{itemize}

\section*{Bemerkenswerte Eigenprojekte}

\begin{itemize}
  \item \textbf{VTT (Virtual Tabletop):} Eigenentwicklung eines Virtual-Tabletop-Systems für Pen-and-Paper-Rollenspiele (vergleichbar mit Foundry/Roll20). Node.js + WebSocket Backend (Server als Single Source of Truth), Canvas Renderer, Zonen-/Oberflächenmechaniken, JSON-basierte Regeldefinitionen
  \item \textbf{Comersant, Multiplayer-Brettspiel:} mehrjähriges Eigenprojekt (Monopoly-ähnliches Online-Brettspiel) mit Angular, Node.js, WebSocket. Architekturentscheidungen über ADRs dokumentiert
  \item \textbf{hochstrasser-immobilien.de:} Pro-bono-Entwicklung einer Webseite inkl. SEO-optimiertem deutschsprachigen Content für einen Immobilienmakler in Ulm/Neu-Ulm
  \item \textbf{senyaak.work:} persönliche Webseite, mehrsprachig (DE/EN/RU)
\end{itemize}

\section*{Sprachen}

\cvkvblock{
\cvkv{Russisch}{Muttersprache}
\cvkv{Deutsch}{verhandlungssicher (C1), seit 2009 in Deutschland; Ausbildung und gesamte Berufslaufbahn auf Deutsch}
\cvkv{Englisch}{sehr gut in Wort und Schrift (B2/C1), regelmäßiger Konsum technischer und allgemeiner Inhalte}
}

\section*{Ausbildung}

\cventry{08/2015 – 07/2018}{Fachinformatiker für Anwendungsentwicklung (IHK)}{saltation GmbH \& Co.~KG, Bielefeld}
\cventry{02/2014 – 04/2015}{Fernstudium Informatik}{ils.de}

\section*{Schulbildung}

\cventry{09/2013 – 06/2015}{CJD-Gymnasium}{Versmold}
\cventry{09/2011 – 06/2013}{Realschulabschluss, Akademie Klausenhof}{Rhede}
\cventry{09/2009 – 06/2011}{Hauptschulabschluss, WRSD am Deutenberg}{VS-Schwenningen}
\cventry{09/2001 – 06/2008}{Allgemeinbildende Schule Nr.~14}{Omsk (Russland)}
